\chapter{Preliminaries}\label{chap:preliminaries}

This chapter fixes the notation and semantic conventions used in the remainder
of the thesis.  \Cref{sec:relations-terms} introduces relations, terms,
positions, substitutions, and contexts.  \Cref{sec:trs-strategies} defines term
rewrite systems and distinguishes full, outermost, and innermost rewriting.
\Cref{sec:ptrs} extends these notions to probabilistic term rewrite systems.
The lifted semantics and rewrite sequence trees are presented in
\cref{sec:probabilistic-derivations}, followed by the definitions of AST, PAST,
and SAST in \cref{sec:probabilistic-termination}.  Finally,
\cref{sec:transformation-background} recalls the deterministic transformation
that forms the starting point of this thesis.  Standard background on term
rewriting can be found in \cite{baader_nipkow_1999}; the probabilistic
conventions follow \cite{avanzini_dallago_yamada_2020,kassing_giesl_2025}.

%-----------------------------------------------------------------------
\section{Relations and Terms}\label{sec:relations-terms}
%-----------------------------------------------------------------------

We begin with notation for relations.  It will be used both for ordinary
rewrite relations and for relations between probability distributions.
Throughout the thesis, $\IN=\{0,1,2,\ldots\}$.

\begin{Definition}{Relations, closures, and normal forms}{relation-notation}
  Let $A$ be a set and let $\to\;\subseteq A\times A$ be a binary relation.
  For $n\in\IN$, its $n$-fold composition $\to^n$ is defined by
  \[
    a\to^0 b \quad\Longleftrightarrow\quad a=b
  \]
  and
  \[
    a\to^{n+1}b
    \quad\Longleftrightarrow\quad
    \text{there is a }c\in A\text{ with }a\to^n c\text{ and }c\to b.
  \]
  The relations
  \[
    \to^*\;=\bigcup_{n\in\IN}\to^n
    \qquad\text{and}\qquad
    \to^+\;=\bigcup_{n\geq 1}\to^n
  \]
  are the reflexive-transitive and transitive closures of $\to$,
  respectively.  An element $a\in A$ is a \emph{normal form} if no $b\in A$
  satisfies $a\to b$.  The set of all normal forms is denoted by
  \[
    \NF_{\to}=\{a\in A\mid \nexists b\in A.\ a\to b\}.
  \]
  The relation $\to$ is \emph{terminating}, or \emph{well-founded in the
  forward direction}, if there is no infinite sequence
  $a_0\to a_1\to a_2\to\cdots$.
\end{Definition}

A term is a finite tree whose internal nodes are function symbols and whose
leaves may additionally be variables.

\begin{Definition}{Signatures and terms}{terms}
  A \emph{signature} $\Sigma=\biguplus_{n\in\IN}\Sigma_n$ is a finite set of
  function symbols with fixed arities; a symbol $f\in\Sigma_n$ has arity
  $n$.  Let $\Var$ be a countably infinite set of variables disjoint from
  $\Sigma$.  The set $\TSet{\Sigma}{\Var}$ of terms over $\Sigma$ and $\Var$
  is the least set satisfying:
  \begin{itemize}
    \item every $x\in\Var$ is in $\TSet{\Sigma}{\Var}$; and
    \item if $f\in\Sigma_n$ and $t_1,\ldots,t_n\in\TSet{\Sigma}{\Var}$,
          then $f(t_1,\ldots,t_n)\in\TSet{\Sigma}{\Var}$.
  \end{itemize}
  Constants are the symbols of arity zero.  A term without variables is
  \emph{ground}; the set of ground terms is denoted by $\TSet{\Sigma}{\emptyset}$.
  For a term $t$, $\Var(t)$ is the set of variables occurring in $t$.  If
  $t=f(t_1,\ldots,t_n)$, then $\rootterm(t)=f$.
\end{Definition}

For example, with binary $\mathsf{pair}$, unary $\mathsf{first}$, and constants
$\mathsf a$ and $\mathsf{produce}$, the expression
$\mathsf{first}(\mathsf{pair}(\mathsf a,\mathsf{produce}))$ is a ground term.

Positions identify nodes in the tree representation of a term.  We use the
empty sequence $\varepsilon$ for the root and concatenate position components
with a dot.

\begin{Definition}{Positions, subterms, and replacement}{positions}
  The set of positions of a term is defined recursively by
  \[
    \operatorname{Pos}(x)=\{\varepsilon\}
  \]
  for variables $x$, and
  \[
    \operatorname{Pos}(f(t_1,\ldots,t_n))
      =\{\varepsilon\}\cup
        \{i.\pi\mid 1\leq i\leq n,\ \pi\in\operatorname{Pos}(t_i)\}.
  \]
  For $\pi\in\operatorname{Pos}(t)$, the subterm of $t$ at $\pi$ is denoted
  by $t|_\pi$.  Replacing this subterm by $s$ yields the term $t[s]_\pi$.

  For positions $\pi$ and $\tau$, we write $\pi\leq\tau$ if $\pi$ is a
  prefix of $\tau$, and $\pi<\tau$ if the prefix is proper.  In this case,
  $\pi$ is \emph{above} $\tau$ and $\tau$ is \emph{below} $\pi$.  Positions
  are \emph{parallel}, written $\pi\parallel\tau$, if neither is a prefix of
  the other.
\end{Definition}

For
$t=\mathsf{first}(\mathsf{pair}(\mathsf a,\mathsf{produce}))$, we have
$t|_{1.2}=\mathsf{produce}$ and
\[
  t[\mathsf b]_{1.2}
  =\mathsf{first}(\mathsf{pair}(\mathsf a,\mathsf b)).
\]
Moreover, $\varepsilon<1<1.2$.

\begin{Definition}{Substitutions and contexts}{substitutions-contexts}
  A \emph{substitution} is a mapping
  $\sigma:\Var\to\TSet{\Sigma}{\Var}$ that differs from the identity on only
  finitely many variables.  It extends homomorphically to terms:
  \[
    f(t_1,\ldots,t_n)\sigma
      =f(t_1\sigma,\ldots,t_n\sigma).
  \]
  We call $t\sigma$ an \emph{instance} of $t$.

  A \emph{context} $C$ is a term containing exactly one occurrence of a fresh
  hole symbol $\Box$.  The term obtained by replacing the hole with $t$ is
  denoted by $C[t]$.
\end{Definition}

Thus, a term $s$ occurs as a subterm of $t$ exactly when $t=C[s]$ for a suitable
context $C$.

%-----------------------------------------------------------------------
\section{Term Rewrite Systems and Strategies}\label{sec:trs-strategies}
%-----------------------------------------------------------------------

\begin{Definition}{Term rewrite systems}{trs}
  A \emph{rewrite rule} is a pair $\ell\to r$ of terms such that
  $\ell\notin\Var$ and $\Var(r)\subseteq\Var(\ell)$.  A \emph{term rewrite
  system} (TRS) $\R$ is a finite set of rewrite rules.

  The root symbols of left-hand sides are the \emph{defined symbols} of
  $\R$.  The remaining symbols of the signature are called
  \emph{constructors}.
\end{Definition}

The variable condition prevents a rewrite step from introducing arbitrary new
terms through previously unbound variables.

\begin{Definition}{Rewrite steps and redexes}{rewrite-step}
  Let $s,t\in\TSet{\Sigma}{\Var}$.  We write
  \[
    s\fto_{\R,\pi}t
  \]
  if there are a rule $\ell\to r\in\R$ and a substitution $\sigma$ such that
  \[
    s|_\pi=\ell\sigma
    \qquad\text{and}\qquad
    t=s[r\sigma]_\pi.
  \]
  The occurrence $s|_\pi=\ell\sigma$ is a \emph{redex}, and $\pi$ is a
  \emph{redex position}.  We omit $\pi$ when its value is irrelevant and
  write $s\fto_\R t$.
\end{Definition}

The next definition restricts which redex positions may be selected.  Stating
the conditions in terms of redex positions, rather than the eventual target
term, is important: a position is disallowed whenever any rule is applicable
strictly above or below it.

\begin{Definition}{Outermost and innermost rewriting}{rewrite-strategies}
  Let $\operatorname{Red}_{\R}(s)$ be the set of redex positions in $s$.
  A step $s\fto_{\R,\pi}t$ is
  \begin{itemize}
    \item \emph{outermost}, written $s\oto_{\R,\pi}t$, if there is no
          $\tau\in\operatorname{Red}_{\R}(s)$ with $\tau<\pi$; and
    \item \emph{innermost}, written $s\ito_{\R,\pi}t$, if there is no
          $\tau\in\operatorname{Red}_{\R}(s)$ with $\pi<\tau$.
  \end{itemize}
  We write $\fto_\R$ for unrestricted, or \emph{full}, rewriting and omit
  positions from $\oto_\R$ and $\ito_\R$ when they are not needed.
\end{Definition}

Equivalently, an innermost step contracts a redex whose proper subterms are all
normal forms.  Every reducible finite term has both an outermost and an
innermost redex.  Consequently, full, outermost, and innermost rewriting have
the same normal forms, although their infinite derivations can differ.

\begin{Definition}{Strategy-specific termination}{strategy-termination}
  A TRS $\R$ is \emph{outermost terminating} if $\oto_\R$ is terminating,
  and it is \emph{innermost terminating} if $\ito_\R$ is terminating.
  Full termination is termination of $\fto_\R$.
\end{Definition}

A term is \emph{linear} if every variable occurs at most once.  A TRS is
left-linear (right-linear) if every left-hand side (right-hand side) is linear,
and it is linear if both conditions hold.  These properties are not required by
the deterministic transformation of \cite{thiemann_2009_outermost}, but they
occur in related strategy-transfer criteria and will be mentioned when those
criteria are compared.

%-----------------------------------------------------------------------
\section{Probabilistic Term Rewrite Systems}\label{sec:ptrs}
%-----------------------------------------------------------------------

Probabilistic rules return finite multi-distributions as in
\cite{avanzini_dallago_yamada_2020,kassing_giesl_2025}.  The prefix
\enquote{multi} records that two alternatives may lead to syntactically equal
terms while still being retained as separate occurrences.  This is convenient
when individual branches are mapped by a simulation proof.

\begin{Definition}{Finite multi-distributions}{finite-multidistributions}
  Let $A$ be a non-empty set.  A \emph{finite multi-distribution} over $A$ is
  a finite multiset
  \[
    \mu=\{p_1:a_1,\ldots,p_k:a_k\}
  \]
  with $a_i\in A$, $0<p_i\leq 1$, and $\sum_{i=1}^k p_i=1$.
  The set of all finite multi-distributions over $A$ is denoted by
  $\FDist(A)$, and $\Supp(\mu)$ denotes the multiset of elements carrying
  positive weight.

  For $0<q\leq1$, scaling is defined by
  \[
    q\cdot\mu=\{qp_1:a_1,\ldots,qp_k:a_k\}.
  \]
  If $F:A\to B$ is a function, its pointwise lifting to multi-distributions is
  \[
    F_\#(\mu)=\{p_1:F(a_1),\ldots,p_k:F(a_k)\}.
  \]
\end{Definition}

The scaled object $q\cdot\mu$ is a weighted multiset of total weight $q$, not
in general an element of $\FDist(A)$.  Such objects are combined with multiset
union so that the resulting weights again sum to one.

\begin{Definition}{Probabilistic term rewrite systems}{ptrs-definition}
  A \emph{probabilistic rewrite rule} has the form
  \[
    \ell\to\mu,
    \qquad
    \mu\in\FDist(\TSet{\Sigma}{\Var}),
  \]
  where $\ell\notin\Var$ and
  $\Var(r)\subseteq\Var(\ell)$ for every $r\in\Supp(\mu)$.  A
  \emph{probabilistic term rewrite system} (PTRS) $\P$ is a finite set of
  probabilistic rewrite rules.  We write a rule explicitly as
  \[
    \ell\to\{p_1:r_1,\ldots,p_k:r_k\}.
  \]
  An ordinary rewrite rule $\ell\to r$ is identified with the deterministic
  probabilistic rule $\ell\to\{1:r\}$.
\end{Definition}

PTRSs may be defined more generally as countable sets of rules.  This thesis
restricts attention to finite PTRSs over finite signatures, as required for the
transformation to produce a finite system.

One step still makes a nondeterministic choice of rule and redex position.  Once
that choice has been made, the right-hand side determines the probabilistic
outcomes.

\begin{Definition}{Probabilistic rewrite steps}{probabilistic-rewrite-step}
  Let $\ell\to\{p_1:r_1,\ldots,p_k:r_k\}\in\P$.  If
  $s|_\pi=\ell\sigma$, then $s$ rewrites at $\pi$ to the multi-distribution
  \[
    \mu=\{p_1:s[r_1\sigma]_\pi,\ldots,
             p_k:s[r_k\sigma]_\pi\}.
  \]
  This unrestricted step is written $s\fto_{\P,\pi}\mu$.  It is written
  $s\oto_{\P,\pi}\mu$ if $\pi$ is an outermost redex position and
  $s\ito_{\P,\pi}\mu$ if $\pi$ is an innermost redex position, using the
  strategy conditions from \cref{sec:trs-strategies}.
\end{Definition}

Normal forms and redex positions depend only on the left-hand sides of $\P$;
the probabilities influence the successors, but not where a rule is
applicable.  We therefore write $\operatorname{Red}_{\P}(s)$ for the positions
at which a left-hand side of $\P$ matches a subterm of $s$.

\begin{example}[Running example]\label{ex:probabilistic-strategy-difference}
  Let $\P_{\mathsf{ex}}$ contain the rules
  \[
    \begin{aligned}
      \mathsf{produce}
        &\to
        \{\tfrac12:\mathsf{pair}(\mathsf a,\mathsf{produce}),
          \tfrac12:\mathsf{pair}(\mathsf b,\mathsf{produce})\},\\
      \mathsf{first}(\mathsf{pair}(x,y))
        &\to\{1:x\}.
    \end{aligned}
  \]
  The first probabilistic step from $\mathsf{first}(\mathsf{produce})$ is both
  outermost and innermost, since the occurrence of $\mathsf{produce}$ is the
  only redex.  It yields
  \[
    \left\{
      \tfrac12:\mathsf{first}(\mathsf{pair}(\mathsf a,\mathsf{produce})),
      \tfrac12:\mathsf{first}(\mathsf{pair}(\mathsf b,\mathsf{produce}))
    \right\}.
  \]
  In each outcome, the root redex is outermost, while the nested
  $\mathsf{produce}$ redex is innermost.  Thus the strategy determines which
  redex must be selected next.
\end{example}

%-----------------------------------------------------------------------
\section{Probabilistic Derivations}\label{sec:probabilistic-derivations}
%-----------------------------------------------------------------------

A term-to-distribution relation describes one probabilistic choice.  For a
relation $\to\;\subseteq A\times\FDist(A)$, we use
\[
  \NF_{\to}=\{a\in A\mid \nexists\mu\in\FDist(A).\ a\to\mu\}
\]
for its normal forms.  To follow
the probability mass through several steps, it is lifted to a relation between
multi-distributions.  Normal forms are retained by a self-loop, so all lifted
derivations can be viewed as infinite sequences.  The following lifting is the
standard multidistribution semantics of probabilistic rewriting
\cite{avanzini_dallago_yamada_2020,kassing_giesl_2025}.

\begin{Definition}{Lifting to multi-distributions}{lifting}
  Let $\to\;\subseteq A\times\FDist(A)$.  Its \emph{lifting}
  $\Rightarrow\;\subseteq\FDist(A)\times\FDist(A)$ is defined as follows.
  For
  \[
    \mu=\{p_1:a_1,\ldots,p_k:a_k\},
  \]
  choose, for every occurrence $a_i$, a multi-distribution $\nu_i$ such that
  \[
    \nu_i=\{1:a_i\}\quad\text{if }a_i\in\NF_{\to},
    \qquad\text{or}\qquad
    a_i\to\nu_i.
  \]
  Then
  \[
    \mu\Rightarrow\biguplus_{i=1}^k p_i\cdot\nu_i.
  \]
  Different legal steps may be chosen for different occurrences of the same
  term.  Thus the lifting retains the nondeterminism of the underlying rewrite
  relation.
\end{Definition}

For a PTRS $\P$, the liftings of $\fto_\P$, $\oto_\P$, and $\ito_\P$ are
written $\fliftto_\P$, $\oliftto_\P$, and $\iliftto_\P$, respectively.

In the running example, outermost rewriting gives the lifted derivation
\[
  \begin{aligned}
    \{1:\mathsf{first}(\mathsf{produce})\}
    &\oliftto_{\P_{\mathsf{ex}}}
      \{\tfrac12:\mathsf{first}(\mathsf{pair}(\mathsf a,\mathsf{produce})),
        \tfrac12:\mathsf{first}(\mathsf{pair}(\mathsf b,\mathsf{produce}))\}
      \\
    &\oliftto_{\P_{\mathsf{ex}}}
      \{\tfrac12:\mathsf a,\tfrac12:\mathsf b\}.
  \end{aligned}
\]
Afterwards, both normal forms are retained.  Under innermost rewriting, the
second lifted step instead expands both nested occurrences of
$\mathsf{produce}$, and this process continues without producing normal-form
probability mass.

Rewrite sequence trees provide an equivalent, and often more convenient,
representation.  Probabilistic alternatives are represented by branches,
whereas the choice of a rule and a redex position is fixed locally at each
internal node \cite{kassing_giesl_2025}.

\begin{Definition}{Rewrite sequence trees}{rewrite-sequence-tree}
  Let $\to\;\subseteq A\times\FDist(A)$.  A \emph{rewrite sequence tree}
  (RST) $\mathfrak T$ for $\to$ is a rooted, finitely branching tree whose
  nodes $v$ are labelled by pairs $(p_v:a_v)$ with $0<p_v\leq1$.
  The root has probability $1$.  If a node $v$ has children
  $w_1,\ldots,w_k$, then
  \[
    a_v\to
    \left\{
      \frac{p_{w_1}}{p_v}:a_{w_1},\ldots,
      \frac{p_{w_k}}{p_v}:a_{w_k}
    \right\}.
  \]
  Leaves are permitted to be reducible.  The tree is \emph{fully evaluated} if
  every leaf is labelled by a normal form.  We write
  $N^{\mathfrak T}$ for the nodes and $\operatorname{Leaf}^{\mathfrak T}$ for
  the leaves.  The \emph{convergence probability} of the tree is
  \[
    |\mathfrak T|
      =\sum_{v\in\operatorname{Leaf}^{\mathfrak T}}p_v.
  \]
\end{Definition}

In a partial RST, reducible leaves form an unfinished frontier, so their weight
is included in $|\mathfrak T|$ even though they have not reached normal form.
For a fully evaluated RST, all leaf weight is normal-form weight and
$|\mathfrak T|$ directly represents the probability of termination.

An infinite branch contributes no leaf probability.  Nevertheless, the total
probability mass that never reaches a leaf can be zero even if infinite paths
exist.  This distinction is the reason why ordinary well-foundedness is too
strong for probabilistic systems.

%-----------------------------------------------------------------------
\section{Probabilistic Termination}\label{sec:probabilistic-termination}
%-----------------------------------------------------------------------

For a multi-distribution, the probability mass that has already reached a
normal form is the following quantity.

\begin{Definition}{Normal-form probability}{normal-form-probability}
  Let $\to\;\subseteq A\times\FDist(A)$ and
  $\mu\in\FDist(A)$.  Its \emph{normal-form probability} is
  \[
    |\mu|_{\to}
      =\sum_{\substack{(p:a)\in\mu\\a\in\NF_{\to}}}p.
  \]
\end{Definition}

Since the lifting retains normal forms, the sequence
$|\mu_0|_{\to},|\mu_1|_{\to},\ldots$ is non-decreasing along a lifted
derivation and bounded by one.  Its limit therefore exists.

\begin{Definition}{Almost-sure termination
  \cite{avanzini_dallago_yamada_2020}}{almost-sure-termination}
  A term $a\in A$ is \emph{almost-surely terminating} with respect to
  $\to$ if every infinite lifted derivation
  \[
    \mu_0\Rightarrow\mu_1\Rightarrow\mu_2\Rightarrow\cdots,
    \qquad \mu_0=\{1:a\},
  \]
  satisfies
  \[
    \lim_{n\to\infty}|\mu_n|_{\to}=1.
  \]
  The relation $\to$ is \emph{almost-surely terminating}, written
  $\AST[\to]$, if every $a\in A$ is almost-surely terminating.

  Equivalently, $\AST[\to]$ holds if every RST for $\to$ has convergence
  probability one; it suffices to consider fully evaluated RSTs.
\end{Definition}

This termwise formulation is equivalent to quantifying over infinite lifted
derivations from arbitrary finite initial multi-distributions, since the
finitely many weighted occurrences can be considered separately.

For a PTRS $\P$, \emph{outermost AST} is $\AST[\oto_\P]$, and
\emph{innermost AST} is $\AST[\ito_\P]$.  The quantification over every lifted
derivation, or every RST, is essential.  Probabilistic choices occur as branches
inside a tree, while different rule or redex choices yield different trees.  AST
requires probability-one termination for all resolutions of this
nondeterminism.

\begin{example}\label{ex:ast-versus-sure-termination}
  The single-rule PTRS
  \[
    \mathsf{retry}
      \to\{\tfrac12:\mathsf{done},\tfrac12:\mathsf{retry}\}
  \]
  has an infinite branch that chooses $\mathsf{retry}$ every time.  After $n$
  steps, however, the probability mass still assigned to $\mathsf{retry}$ is
  $2^{-n}$.  Hence the normal-form probability tends to one, and the system is
  AST even though it is not terminating in the classical sense.
\end{example}

In \cref{ex:probabilistic-strategy-difference}, the term
$\mathsf{first}(\mathsf{produce})$ is almost-surely terminating under outermost
rewriting: all probability mass reaches $\mathsf a$ or $\mathsf b$ after two
steps.  Under innermost rewriting, its normal-form probability remains zero.
This is the strategy gap that the transformation in \cref{chap:main-chapter}
is designed to bridge.

Two stronger probabilistic termination notions are useful for distinguishing
probability-one termination from finite expected runtime.

\begin{Definition}{PAST and SAST
  \cite{avanzini_dallago_yamada_2020,kassing_giesl_2025}}{past-sast}
  Let $\vec\mu=(\mu_n)_{n\in\IN}$ be an infinite lifted derivation.  Its
  \emph{expected derivation length} is
  \[
    \edl(\vec\mu)
      =\sum_{n=0}^{\infty}\bigl(1-|\mu_n|_{\to}\bigr).
  \]
  The relation $\to$ is \emph{positively almost-surely terminating}, written
  $\PAST[\to]$, if $\edl(\vec\mu)<\infty$ for every infinite lifted derivation
  with $\mu_0=\{1:a\}$ for some $a\in A$.

  For $a\in A$, its \emph{expected derivation height} is
  \[
    \edh_{\to}(a)
      =\sup\{\edl(\vec\mu)\mid
        \vec\mu=(\mu_n)_{n\in\IN}\text{ is a lifted derivation and }
        \mu_0=\{1:a\}\}.
  \]
  The relation is \emph{strongly almost-surely terminating}, written
  $\SAST[\to]$, if $\edh_{\to}(a)<\infty$ for every $a\in A$.
\end{Definition}

The notions form the strict hierarchy
\[
  \SAST[\to]\quad\Longrightarrow\quad
  \PAST[\to]\quad\Longrightarrow\quad
  \AST[\to],
\]
in general \cite{avanzini_dallago_yamada_2020,kassing_giesl_2025}.  This thesis
focuses on AST; PAST and SAST are included to make the chosen property precise
and to prevent it from being confused with finite expected derivation length.
PAST requires each nondeterministic derivation to have finite expected length;
SAST additionally requires a uniform finite bound for all derivations from each
fixed start term.

%-----------------------------------------------------------------------
\section{Transformation-Based Strategy Transfer}
\label{sec:transformation-background}
%-----------------------------------------------------------------------

The transformation in this thesis is based on the deterministic construction
of Thiemann \cite{thiemann_2009_outermost}.  Its extended signature contains
fresh control symbols.  Starting from a marked source term, the transformed
system first moves a search marker towards a redex.  Innermost rewriting permits
the marker to descend only when the surrounding source term is not itself a
redex; this enforces the outermost restriction.  A transformed source rule then
performs the actual rewrite step and changes the marker into a return marker.
Finally, deterministic administrative rules move that marker back to the root
before the next simulated step begins.

\begin{Theorem}{Deterministic outermost-to-innermost transformation
  \cite{thiemann_2009_outermost}}{deterministic-transformation}
  For every finite TRS $\R$ over a finite signature, the transformation
  constructs a finite TRS $\R'$ such that
  \[
    \R\text{ is outermost terminating}
    \quad\Longleftrightarrow\quad
    \R'\text{ is innermost terminating}.
  \]
\end{Theorem}

\paragraph{Proof idea.}
The forward simulation maps every outermost source step to a non-empty finite
innermost derivation.  For the converse, arbitrary infinite innermost
derivations over the extended signature are analysed and a suitable derivation
is projected back to an infinite outermost source derivation.  The complete
proof is given in \cite{thiemann_2009_outermost}.

An earlier transformation by Raffelsieper and Zantema instead reduces
outermost ground termination to unrestricted termination.  Its principal
correctness result is one-way: termination of the transformed system implies
outermost ground termination of the source.  For finite quasi-left-linear TRSs,
the required anti-matching information and the transformed system can be
constructed finitely \cite{raffelsieper_zantema_2009}.  This result is useful as
a second example of the simulation-based proof pattern, but it must not be
confused with the outermost-to-innermost equivalence in
\cref{thm:deterministic-transformation}.

These deterministic simulation results motivate the probabilistic construction,
but they do not by themselves establish an AST result.  The new proof must
track probability mass through infinite lifted derivations or fully evaluated
RSTs and must represent every nondeterministic outermost source choice.  The concrete
encoding, the transformed PTRS, and these correctness arguments are introduced
in \cref{chap:main-chapter}.
